% Seção 7 — Ameaças à identificação. Flags 0–7 do CLAUDE.md + docs/ars/19 §6.

\chapter{Ameaças à identificação}
\label{cap:ameacas}

O Quadro~\ref{qua:ameacas} reúne as ameaças conhecidas, o que cada uma faz ao
parâmetro estimado e o que o desenho lhe opõe. Três merecem comentário.

\begin{quadro}[htb]
\caption{\label{qua:ameacas}Ameaças à identificação e respostas do desenho}
\centering
\footnotesize
\begin{tabular}{@{}p{3.5cm}p{4.9cm}p{6.1cm}@{}}
\toprule
Ameaça & Efeito sobre a estimativa & Resposta \\
\midrule
Tratamento medido pela cultura, não pelo método &
Atenuação proporcional ao $VPP$ (Eq.~\ref{eq:atenuacao}) &
Estágio 2 (medição) e estágio 3 (tratados pelo método). \\
Prática encerrada antes da lei (2011) &
Grupo tratado inexistente; nulo por ausência de tratamento &
Autos judiciais e relatórios do MAPA; regra de decisão (Quadro~\ref{qua:regra}). \\
Antecipação (2011, 2015 e 2018) &
Dose e desfecho respondem antes da vigência &
Sensibilidade com dose de 2010--2014; coortes de transição excluídas; \emph{leads} até 2018; o estágio 3 não usa a dose. \\
Grupo de comparação pequeno e contaminado &
Nível da estimativa depende de quem é o zero; controle vetorial aéreo também foi proibido &
Quatro construções do grupo; inferência adequada a poucos grupos; estágio 3 não depende desse grupo. \\
Seleção para o nascimento vivo &
Efeito sobre o peso médio atenuado ou invertido &
Óbito fetal como desfecho confirmatório. \\
Composição sazonal das mães &
Hiato sazonal reflete quem concebe em cada época &
Efeitos fixos de município por mês do ano; teste de estabilidade da composição materna. \\
Choques de 2020--2021 &
Mudança de sazonalidade de concepção e de assistência pré-natal &
Estimativas com e sem o período; doadores externos. \\
Substituição pelo aplicador terrestre &
Exposição populacional cai, ocupacional pode subir &
Notificações de intoxicação como canal; faz parte do efeito da política. \\
Fiscalização desconhecida &
Parâmetro é de intenção de tratar &
Declarado; registros de autuação solicitados nos pedidos de acesso. \\
\bottomrule
\end{tabular}
\fonte{Elaboração própria.}
\end{quadro}

\textbf{A medida do tratamento.} É a ameaça que organiza o projeto. A
proibição alcança o avião, e a variável disponível mede a lavoura. Se a
aplicação aérea era concentrada em poucos municípios, como indica o Censo
Agropecuário de 2006, a maior parte dos
classificados como tratados pela cultura não o era pelo método, e a
Equação~\ref{eq:atenuacao} implica que a forma reduzida pode ser pequena mesmo
que o efeito verdadeiro seja grande. Ampliar o grupo tratado pela cultura
tenderia a agravar o problema, ao acrescentar municípios sem registro de
avião. A resposta é medir o método e concentrar a
comparação onde ele existia.

\textbf{A persistência da prática.} Todas as estratégias que medem o método
dependem de a pulverização aérea ter continuado na Chapada do Apodi entre 2011
e 2018. O fato de a lei de 2019 ter sido aprovada contra uma prática descrita
como corrente sugere que sim, mas não o demonstra. Se os registros mostrarem
que ela cessou, o resultado do ensaio muda de natureza: a ausência de efeito
passa a ser ausência de tratamento, e a contribuição se desloca para a
mensuração do alcance efetivo da lei.

\textbf{A sazonalidade.} O desenho de calendário troca a hipótese de
tendências paralelas entre municípios pela de estabilidade relativa do hiato
sazonal. A troca é vantajosa se a aplicação aérea tiver de fato um pico
sazonal; se a aplicação for uniforme ao longo do ano, o desenho mede zero por
construção. Por isso a existência do pico é verificada nos relatórios mensais
antes da estimação, e o poder do desenho é calculado no pré-período.
